\chapter{排除犯罪事由II：排除责任事由}

% Source PDF physical page 107; see audit/source_coverage.json.
\begin{EditionOpeningPanel}
1.行为人制造了违法事实后，接下来的审查任务是，就该违法事实而言，能否谴责行为人？行为人是否具有可谴责性？这里的谴责被称为“责任”。正常情况下，只要行为人制造了违法事实，就能谴责行为人。但是，如果行为人具有排除责任（谴责）的事由（也称为宽恕事由、责任阻却事由），则无法谴责行为人，行为人最终不用承担刑事责任（不利的法律后果），最终无罪。

2.重点：责任年龄。难点：法律认识错误。
\end{EditionOpeningPanel}

\section{责任年龄}

[问题]15周岁的狗蛋闯荡江湖，干起了绑票的生意，将小芳绑架了，由于小芳丈夫不给钱，狗蛋就撕票。对狗蛋如何处理？（绑票案，2020年试题）

\InsertDiagram{figures/ch10_age.tex}

责任年龄，是指刑法所规定的，行为人能够接受刑法谴责所须达到的年龄。如果没有达到责任年龄，即使行为人制造了法益侵害事实，主观上具有故意或过失，也无法谴责行为人，进而不用负刑事责任。责任年龄分为三个阶段。

\subsection*{一、完全无责任年龄}

这是指不满12周岁的人。不满12周岁的人对所有犯罪都不用负刑事责任。

[注意1]生日的第二天才算满1周岁（2016年第3题）。例如，甲1990年8月8日出生，2002年8月8日生日这天仍属于不满12周岁，第二天（2002年8月9日）才算已满12周岁。

[注意2]犯罪嫌疑人不讲自己的年龄，司法机关可以进行骨龄鉴定，鉴定意见证明犯罪嫌疑人已满12周岁，就应负相应的刑事责任。这种鉴定不要求查明具体的出生日期，如几月几日。（2016年第3题）

% Source PDF physical page 108; see audit/source_coverage.json.
\subsection*{二、相对责任年龄}

\subsubsection*{（一）12 至14 周岁}

第17条第3款已满十二周岁不满十四周岁的人，犯故意杀人、故意伤害罪，致人死亡或者以特别残忍手段致人重伤造成严重残疾，情节恶劣，经最高人民检察院核准追诉的，应当负刑事责任。【《刑法修正案（十一）》增设］

1.这两种罪是指犯罪行为，不限于两种罪名。例如，在抢劫中故意杀人的，在强奸中故意杀人的，在绑架中故意杀人的，均属于第17条第3款中的故意杀人。

2.这两种罪要求造成实害结果。注意，重伤与严重残疾是并且的关系。

3.这两种罪都要求情节恶劣。\footnote{参见2026年7月31日最高人民检察院、公安部《关于办理已满十二周岁不满十四周岁未成年人严重暴力犯罪核准追诉案件若干问题的规定》。}

4.程序上，要求经最高人民检察院核准追诉。

\subsubsection*{（二）14至16 周岁}

第17条第2款已满十四周岁不满十六周岁的人，犯\CJKunderwave{故意杀人}、\CJKunderwave{故意伤害致人重伤}\CJKunderwave{或者死亡}、\CJKunderwave{强奸}、\CJKunderwave{抢劫}、\CJKunderwave{贩卖毒品}、\CJKunderwave{放火}、\CJKunderwave{爆炸}、\CJKunderwave{投放危险物质}罪的，应当负刑事责任。“烧杀淫掠，商贩暴徒（伤贩爆投）”

1.这八种罪都是严重犯罪，都是故意犯罪。

[注意1]八种罪里不包括：失火罪，绑架罪，拐卖妇女、儿童罪，决水罪，走私、制造、运输毒品罪。考试常用一个不负责任的行为附带一个要负责任的行为。

例1，本节开头“绑票案”中，15周岁的狗蛋对绑架罪不负责任，但对故意杀人罪要负责任。（2020年试题）

例2，15周岁的狗蛋拐卖妇女并强奸妇女，对拐卖妇女罪不负责任，但对强奸罪要负责任。

例3，15周岁的狗蛋为强迫妇女卖淫，先强奸妇女，然后强迫其卖淫，对强迫卖淫罪不负责任，但对强奸罪要负责任。

例4，15周岁的狗蛋实施妨害公务的行为，并将人打成重伤，对妨害公务罪不负责任，但对故意伤害罪（重伤）要负责任。

[注意2］故意伤害致人重伤或者死亡，其中：（1）“致人死亡”是指过失致人死亡。

（2）“致人重伤”是指故意致人重伤。例如，15周岁的狗蛋只想故意致人轻伤，但过失致人重伤，不负刑事责任。（2026年试题）

2.这八种罪的范围：这八种罪是指八种犯罪行为，而不仅限于八个罪名。

（1）这八种罪包括法条竞合来的八种罪。例如，包括抢劫枪支、弹药、爆炸物、危险物质罪。（2026年试题）

（2）这八种罪包括法律拟制来的八种罪。所谓法律拟制，是指原本不符合此罪要件，但法律硬是按此罪论处，特事特办。\footnote{关于法律拟制，具体参见第十七讲中的“注意规定与法律拟制”。}

第一，法律拟制来的故意杀人罪，共有5个。实施以下5个犯罪过程中，过失致人死亡，

% Source PDF physical page 109; see audit/source_coverage.json.
定故意杀人罪：非法拘禁罪，聚众斗殴罪，刑讯逼供罪，暴力取证罪，虐待被监管人罪。

例如，狗蛋15周岁时，是个古惑仔，参与聚众斗殴，过失致人死亡。对狗蛋定故意杀人罪。

第二，法律拟制来的抢劫罪，共有3个：第267条第2款（携带凶器抢夺定抢劫），第269条（事后转化抢劫），第289条（聚众“打砸抢”中毁坏、抢夺财物定抢劫）。

例如，狗蛋15周岁时，携带凶器抢夺小芳提包。对狗蛋定抢劫罪。（2009年第2题）

[例外］事后转化抢劫。司法解释规定：14至16周岁的人，先实施盗窃、诈骗、抢夺，为了窝藏赃物、抗拒抓捕、毁灭罪证，当场使用暴力，不按照事后转化抢劫罪来处理；如果暴力行为构成故意伤害罪（重伤）、故意杀人罪，则按照这两罪处理。该司法解释属于特殊规定，需要作为例外来记。（2010年第4题）

3.14至16周岁的人实施八种罪的帮助行为，不负刑事责任。（2009年第2题，2010年第4题)

\subsection*{三、完全责任年龄}

第17条第1款已满十六周岁的人犯罪，应当负刑事责任。

第17条第5款因不满十六周岁不予刑事处罚的，责令其父母或者其他监护人加以管教；在必要的时候，\CJKunderwave{依法进行专门矫治教育}。

[提示］《刑法修正案（十一）》将原来的“收容教养”改为“矫治教育”。

\subsection*{四、减轻责任年龄}

1.12至18周岁。第17条第4款规定，已满12周岁不满18周岁的人犯罪，应当从轻或者减轻处罚。

2.已满75周岁。第17条之一规定：已满75周岁的人\CJKunderwave{故意犯罪的}，\CJKunderwave{可以}从轻或者减轻处罚；\CJKunderwave{过失犯罪的}，\CJKunderwave{应当}从轻或者减轻处罚。

\section{责任能力}

[问题］由于小芳甩了狗蛋，狗蛋欲杀之，捅了一刀，小芳鲜血直流，狗蛋看到这么多血，很吃惊，导致间歇性精神病发作，在此状态下又捅了一刀。小芳死亡。对狗蛋该如何处理？（见血崩溃案，2016年第3题）

责任能力，是指行为人对自己行为性质的辨认能力与控制能力。行为人需要同时具备辨认能力和控制能力，才有责任能力。如果行为人不具备辨认能力或控制能力，即使其制造了法益侵害事实，也无法谴责他，进而不承担刑事责任。（2023年试题）

\subsection*{一、责任能力的程度}

\InsertDiagram{figures/manual_table_p109_1.tex}

% Source PDF physical page 110; see audit/source_coverage.json.
\ContinuedTable{figures/manual_table_p110_1.tex}

[注意1］间歇性精神病人，在精神正常时，属于完全有刑事责任能力，应当负刑事责任。法条在此没有规定“可以从宽处罚”（第18条第2款）。在丧失责任能力时，不负刑事责任。（2020年试题）

[注意2］心理疾病不等于精神病。例如，抑郁症不属于精神病，抑郁症患者具有完全责任能力。（2011年第4题）

\subsection*{二、特殊人群}

1.又聋又哑的人和盲人

残疾人不等于精神病人。聋哑人、盲人，虽然视觉、听觉、口语能力丧失，但是精神正常，属于完全有责任能力，应当负刑事责任（2017年第3题）。但是考虑到是残疾人，可以从轻、减轻或者免除处罚。\CJKunderwave{注意}：聋哑人必须是又聋又哑才可以减免处罚。

2.醉酒的人

第一种，生理性醉酒，即日常生活中的醉酒，属于完全有刑事责任能力，应当负刑事责任（第18条第4款）。\CJKunderwave{注意}：没有可以从轻、减轻处罚的规定。（2017年第3题）

第二种，病理性醉酒，是指因酒精中毒导致幻觉、妄想等精神病症状，是精神病的一种。这属于完全无刑事责任能力，不负刑事责任。

3.吸毒的人

我国刑法没有将吸毒状态认定为丧失或减轻责任能力的情形，因此吸毒状态仍然被认为完全有责任能力。

（1）吸毒后第一次产生幻觉，在幻觉中实施犯罪，由于没有犯罪故意，可认定为过失犯罪。例如，甲毒瘾很大，有次在家里吸毒，第一次产生幻觉，在幻觉中认为乙要来杀自己，为了防止乙的追杀，将窗帘点燃，酿成火灾，导致邻居房屋被烧毁。甲不构成放火罪，可认定为失火罪。

［提示」在此认可吸毒状态的影响，只是认可其影响了犯罪故意，并不是认可其影响了责任能力。也即甲有责任能力，需要承担刑事责任，只是不构成故意犯罪，而构成过失犯罪。（2016年第3题）

（2）明知自己吸毒后会产生幻觉，利用这种特点实施犯罪，应认定为故意犯罪。

\subsection*{三、行为与责任同时存在原则}

该原则是指，行为人在实施行为时，应同时具有责任要件，才能谴责该行为。常考的责任要件有责任年龄、责任能力。

\subsubsection*{（一）行为与责任年龄同时存在原则}

该原则中的“行为”既包括作为，也包括不作为。（2016年第3题）

例1，甲为了炸毁他人房屋，安装炸弹时只有13周岁，炸弹爆炸时已经14周岁。甲

% Source PDF physical page 111; see audit/source_coverage.json.
在14周岁时对其13周岁时安装的炸弹负有拆除义务，也即具有防止结果发生的义务。甲不履行该义务，构成不作为的爆炸罪。（2015年第2题）

例2，甲在12周岁生日当天投毒杀乙，第二天，甲见死不救，构成不作为的故意杀人罪。

\subsubsection*{（二）行为与责任能力同时存在原则}

1.“同时”的要求。仅要求行为与责任能力同时存在，不要求结果与责任能力同时存在。也即，只要求实施行为时有责任能力，不要求发生结果时有责任能力。

例如，甲欲杀乙，邮寄毒酒时甲的精神正常，乙收到毒酒时，甲精神病发作，在此期间，乙喝酒中毒死亡。甲构成故意杀人罪既遂。（2015年第2题）

2.既遂条件。具有责任能力时的犯罪行为是否既遂，关键看该行为与最终结果有无因果关系。

例如，在本节开头“见血崩溃案”中，狗蛋在精神正常时杀人，此时成立故意杀人罪。此时的杀人行为与精神病状态下的杀人行为，与死亡结果属于二因一果。由于精神正常时的杀人行为与死亡结果有因果关系，因此狗蛋构成故意杀人罪既遂。狗蛋在精神病状态下的杀人行为，由于丧失责任能力，不具有可谴责性。（2016年第3题）

[巩固练习]题1，甲欲杀乙，仅砍伤乙的手指后，甲精神病发作，丧失责任能力，又砍断乙的脖子，乙死亡。由于甲在具有责任能力时的行为不具有致命危险，与乙的死亡没有因果关系，因此甲仅构成故意杀人罪未遂。

题2，甲欲杀乙，用刀砍伤乙后，甲精神病发作，丧失责任能力，抢走了乙的财物。甲构成故意杀人罪未遂，对于抢劫行为不负刑事责任。（2015年第2题）

\subsubsection*{（三）例外情形：原因自由行为}

这是指，具有责任能力的人，故意或过失使自己陷人丧失责任能力的状态或无意识的状态，并在该状态下实施了法益侵害行为。

例如，甲明知自已有病理性醉酒史，饮酒后会引发病理性醉酒，丧失责任能力。甲为了伤害乙，故意在乙面前喝醉，丧失责任能力，然后伤害了乙。甲应否承担刑事责任？

让甲承担责任的障碍在于：甲在有责任能力时没有犯罪行为，甲有犯罪行为时又没有责任能力，行为与责任能力不配套，不是同时具备。

但是，主流观点认为，甲应承担刑事责任，认为这是行为与责任同时存在原则的一种

例外情形。让甲承担责任的主要理由是：甲在实施原因行为（导致丧失责任能力的行为）时，具有意志自由、选择自由，甲支配了自己的这种无责任能力的状态。甲构成故意伤害罪既遂。同时\CJKunderwave{注意}：\CJKunderwave{原因自由行为的着手}，不是实施原因行为（如饮酒）时，而是实施结果行为（如伤害行为）时。

[巩固练习]题1，甲欲强奸乙，故意使自已陷入无责任能力状态，然后却实施了抢劫，抢走乙的包。甲的强奸构成犯罪预备。甲在实施抢劫时无责任能力，不负刑事责任。如果甲一开始想抢劫，故意使自已陷入无责任能力状态，然后实施了抢劫，则构成抢劫罪既遂。

题2，甲欲强奸乙，故意使自已陷入无责任能力状态，然后却认错人，强奸了丙。甲构成强奸罪既遂，因为这是同一构成要件内的对象错误，这种对象错误不重要，不影响既遂的认定。

% Source PDF physical page 112; see audit/source_coverage.json.
\section{法律认识错误（违法性认识可能性）}

[问题］狗蛋以为某种蟾是普通动物，便大量猎捕。实际上，该动物是珍贵保护动物。对狗蛋如何处理？（抓某种蟾蜍案）

\subsection*{一、两种认识错误的区分标准}

\subsubsection*{（一）事实认识错误与法律认识错误的区分}

一个人干一件事，通常有两个层次的认知活动：首先，认识到自己在干什么事；然后，认识到自已干的事在刑法上是否被禁止。前者是事实判断，是认识事实。后者是法律评价，是认识刑法的禁止性。前者有可能导致事实认识错误，后者有可能导致法律认识错误。

[结论】事实认识错误，是针对事实产生认识错误。法律认识错误，是针对\CJKunderwave{刑法的禁}\CJKunderwave{止性}产生认识错误。

第一，这里的事实，是指构成要件事实，包括危害行为、行为对象、危害结果。

第二，这里的刑法禁止性，是指刑法对行为人干的事的评价态度：予以禁止。一个行为具有刑法的禁止性，便具有了刑法的违法性。所以，法律认识错误也称为违法性认识错误、禁止性认识错误。\footnote{提示：违法性认识错误中的“违法性”是指刑法的禁止性。排除违法事由中的“违法性”是指法益侵害性。前者被称为形式的违法性，后者被称为实质的违法性。}

例1，甲以为自已贩卖的是普通物品，实际是淫秽物品。甲对行为对象这一事实产生认识错误，属于事实认识错误。

例2（2015年第55题），甲拘禁吸毒的陈某数日。甲认识到自己剥夺了陈某的自由，但误以为刑法不禁止自已的这种行为。甲对自已干了什么事，没有认识错误，但对自己所干的事，刑法是什么评价态度，是禁止还是允许，产生了认识错误，属于法律认识错误。

\subsubsection*{（二）法定犯中的区分问题}

法定犯，是指需要援引某项行政法规范的罪名，也称为行政犯，例如，非法持有枪支罪、危害珍贵野生动物罪、非法经营罪等。自然犯，是指不需要援引行政法规范，古今中外都会规定为犯罪的罪名，例如，放火罪、故意杀人罪、抢劫罪等。事实认识错误与法律认识错误的区分问题主要出现在法定犯中。

第一，\CJKunderwave{法律认识错误中的“法律”，是指刑法（刑法的禁止性）}，不包括行政法。对行政法的认识错误，不是刑法中的法律认识错误。

第二，某项行政法规范是用来解释某个构成要件事实时，对该行政法有认识错误，导致对构成要件事实有认识错误，属于事实认识错误。

例1，在本节开头的“抓某种蟾蜍案”中，危害珍贵野生动物罪的行为对象是“珍贵的国家保护动物”。狗蛋以为抓捕的是普通动物。实际上，《野生动物保护法》将其规定为珍贵的国家保护动物。狗蛋由于对该行政法有认识错误，导致对构成要件事实（行为对象）产生认识错误，没有认识到该动物是“珍贵的国家保护动物”，属于事实认识错误，表明狗蛋不具有实施危害珍贵野生动物罪的犯罪故意。不能因为狗蛋对行政法有认识错误，就认为其存在刑法上的法律认识错误。

例2，非法持有枪支罪的行为对象是“法律禁止私人持有的枪支”。天津大妈赵春华在集市摆了个射击气球的游戏摊位。赵春华认为用于射击气球的游戏枪不是法律禁止持有的真枪。而按照《枪支管理法》规定，该枪属于真枪。赵春华由于对该行政法有认识错

% Source PDF physical page 113; see audit/source_coverage.json.
误，导致对构成要件事实（行为对象）产生认识错误，属于事实认识错误，表明赵春华不具有实施非法持有枪支罪的犯罪故意，不构成犯罪。\footnote{提示：虽然当地法院认定赵春华构成非法持有枪支罪，但这-一判决引起舆论哗然。理论界和实务界主流观点认为，赵春华不构成非法持有枪支罪。}在此不能因为赵春华对行政法有认识错误，就认为其存在刑法上的法律认识错误。

\subsubsection*{（三）两种认识错误的审查顺序}

事实认识错误处在前一板块“主观要件”中。法律认识错误处在后一板块“排除责任事由”中。

第一，在审查顺序上，应先审查事实认识错误，后审查法律认识错误。一旦基于事实认识错误而排除犯罪，就不需要再考虑法律认识错误。

第二，事实认识错误也会引发法律认识错误，但由于在审查体系上，应先审查事实认识错误，审查了事实认识错误，便不需要再考虑此后的法律认识错误。换言之，真正的法律认识错误，前提是没有事实认识错误。

例如，狗蛋误以为自已抓捕的是普通动物，实际是珍贵动物，存在事实认识错误，由此导致其会认为自已的行为不具有刑法禁止性，不构成犯罪，由此产生法律认识错误。但是，由于狗蛋已经构成事实认识错误，按照事实认识错误已经排除了犯罪故意，得出无罪结论，就不需要再审查后续的法律认识错误。换言之，这种情形应归入事实认识错误的范畴，不应按照法律认识错误来处理。许多同学认为，对狗蛋应按照法律认识错误来处理。这种误解的原因在于，没有正确认识两种认识错误的体系地位和审查顺序。

\subsection*{二、两种认识错误的法律后果}

\subsubsection*{（一）事实认识错误}

[结论］对构成要件事实有认识错误，排除犯罪故意，进而不成立故意犯罪。俗语讲：“\CJKunderwave{不知者不为罪}”。其中的“不知”是指不知构成要件事实。

例如，狗蛋没有认识到所抓捕的某种蟾蜍是珍贵保护动物，表明其没有实施危害珍贵野生动物罪的犯罪故意，进而不构成犯罪。

又如，赵春华没有认识到所持有的枪支是真枪，表明其没有实施非法持有枪支罪的犯罪故意，进而不构成犯罪。\footnote{提示：与法律认识错误进行区分的事实认识错误，属于缺乏犯罪故意型事实认识错误。第八讲中的“事实认识错误”是指具有犯罪故意型事实认识错误。}

\subsubsection*{（二）法律认识错误}

案件走到审查法律认识错误这一步，前提是，已经审查了前一板块“主观要件中的故意”，结论是，存在故意，不存在事实认识错误。

事实认识错误是事实判断层面的事情，法律认识错误是法律评价层面的事情。一个人成立故意犯罪，只需要认识到构成要件事实，不需要认识到法律评价，不需要认识到刑法的禁止性（违法性）。法律评价的事情以法官的评价为准，行为人的认识错误不重要，还是能够谴责行为人。俗语讲“\CJKunderwave{不知法者不免责}”便是这个意思。\footnote{例如（2011年第5题），“成立故意犯罪，不要求行为人认识到自己行为的违法性。”该说法是否正确？答案；正确。}

但是，如果行为人连“知法的可能性”都没有，此时还谴责行为人，就说不过去了。此时，“知法的可能性”或“违法性认识的可能性”，便成为排除责任（可谴责性）的一种事由。

% Source PDF physical page 114; see audit/source_coverage.json.
[结论】对行为的违法性（刑法禁止性）有认识错误，有可能排除责任。能否排除责任的条件是：有无违法性认识的可能性。

1.自然犯中，出现法律认识错误，一般都具有违法性认识的可能性。

例1（2002年第4题），甲误以为刑法不禁止征得幼女同意的性交行为，明知乙女只有13周岁，在征得乙的同意后与其发生性交行为。虽然甲没有认识到该行为的违法性（刑法禁止性），存在法律认识错误，但是，没有认识到违法性不等于没有认识到的可能性。奸淫幼女是自然犯，对于甲而言，存在违法性认识的可能性。造成这种法律认识错误是自已的责任，法官可以谴责甲，对甲应以强奸罪论处。

例2（2015年第55题），甲拘禁吸毒的陈某数日。甲误以为刑法不禁止自己的这种行为。非法拘禁罪是自然犯。甲虽然没有认识到行为的违法性（刑法禁止性），但是具有违法性认识的可能性。因此，这种认识错误不排除责任。甲构成非法拘禁罪。

2.法定犯中，出现法律认识错误，有可能缺乏违法性认识的可能性。

缺乏违法性认识可能性的原因：听信有权机关的正式答复。注意：听信一般公民（包括律师、法学专家）的答复，不算数。

例如（2008年第4题），甲在生产经营过程中，对自己的经营活动是否违反刑法，有疑虑，向法院书面咨询其经营是否违反刑法。法院答复：不违反刑法。实际上，甲的经营活动违反刑法。由于甲信赖法院的正式答复，导致法律认识错误，甲没有认识到行为的刑法禁止性（违法性），而且连认识的可能性都没有。此时，法官便无法谴责甲，否则太不公平，对甲不应以犯罪（非法经营罪）论处。

[巩固练习]题1（2015年第55题），“甲知道自已的行为有害，但不知是否违反刑法，遂请教中学语文教师乙，被告知不违法后，甲实施了该行为。但事实上刑法禁止该行为。乙的回答不影响甲成立故意犯罪。”该说法是否正确？答案：正确。若请教体育老师，更不行。（不是歧视体育老师哈。希望听了我的课，作了我的亲学生，以后不要被骂“你的刑法是不是体育老师教的？”）

题2（2025年试题），甲购买了一套别墅，在处理花园草坪时挖出一包毒品。甲认为可能是前房主埋藏的，与自已无关，故将这包毒品又埋回原处，认为自已的行为无罪。第一，非法持有毒品罪是作为犯罪。第二，甲对自已控制持有毒品的事实没有认识错误。第三，甲有法律认识错误，但甲有违法性认识可能性，因此构成犯罪。

[提示]上述法律认识错误是指，行为人误以为自己的行为没有违法性，实际上有。还有种相反情形，行为人误以为自己的行为具有违法性，实际上没有。这种情形被称为幻觉犯（幻想犯），不构成犯罪。对此当然只能按照无罪处理。例如，误以为自已跟人通奸是犯罪，实际上无罪。

\InsertDiagram{figures/manual_p114_1.tex}

% Source PDF physical page 115; see audit/source_coverage.json.
\InsertDiagram{figures/manual_table_p115_1.tex}

\section{期待可能性}

[问题]16周岁的狗蛋偷了他人一万元现金。狗蛋的父亲老狗蛋明知该事实，仍将狗蛋藏起来。对老狗蛋如何处理？（窝藏儿子案）

期待可能性，是指从行为时的具体情况看，可以期待行为人作出合法行为。\CJKunderwave{法律不强}\CJKunderwave{人所难}，如果从行为时的具体情况看，不能期待行为人作出合法行为，行为人即使作出了违法犯罪行为，也不值得谴责。因此，不具有期待可能性便成为排除责任事由。例如，甲杀人后，为了毁灭证据，指使乙帮助碎尸。乙构成帮助毁灭证据罪。但甲不构成该罪的教唆犯，因为法律不能期待犯罪人犯罪后保全证据。

常考情形：

1.近亲属间的窝藏、包庇行为，可以不追究窝藏、包庇罪的刑事责任或应从宽处罚。

例如，本节开头“窝藏儿子案”中，法律无法期待一位父亲出卖自已的儿子，对这位父亲不具有期待可能性，因此可以不追究这位父亲窝藏罪的责任。其实我国汉代就有“亲亲得相首匿”的规定。该规定的古代依据是儒家理论，放到现代，依据就是期待可能性原理。

2.盗窃者将赃物拿到地下销赃市场去卖，购买者也知道是赃物。对这种销赃行为不具有期待可能性，对盗窃者无须定掩饰、隐瞒犯罪所得罪。

\OriginalPanel{典型真题}

间歇性精神病人在不能辨认或者不能控制自已行为时，实施严重危害社会行为的，该如何处理？（2020年回忆版金题）\footnote{[答案]关于间歇性精神病人的刑事责任，根据《刑法》第18条第2款规定，（1）在精神正常时，属于完全有责任能力，应当负刑事责任。法条在此没有规定“可以从宽处罚”。（2）在丧失责任能力时，不负刑事责任。本题答案为B。}

A.应当负刑事责任

B.不负刑事责任

C.应当负刑事责任，但可以减轻或免除处罚

D.待治愈后再追究刑事责任

